% !TEX root = ../main.tex
\addcontentsline{toc}{chapter}{Introduction}
\chapter*{Introduction}

The idea of generating mathematical reasoning computationally
is one of the most profound challenges in computer science.
At the core of this pursuit lies the Curry-Howard correspondence,
an isomorphism between logic and computation \cite{howard}. Under this paradigm,
mathematical propositions are interpreted as types, and the proofs
of propositions are viewed as terms (or programs) inhabiting
those types. Consequently, the problem of proving a theorem is 
reduced to the problem of finding a well-typed term in a formal type-theory.

Modern proof assistants, such as Rocq (formerly Coq) \cite{coquand1988} and Lean \cite{demoura2015},
leverage this correspondence using rich type theories like the 
Calculus of Inductive Constructions \cite{paulin-mohring1993}. While these environments
enabled us to formalize complex theories, fully autonomous theorem
proving with such rich frameworks remains an incredibly
difficult challenge. The primary obstacle is undecidability of 
type inhabitation in dependent type theories: attempting
to exhaustively construct a term of a given type results in a 
combinatorial explosion of the search space.

Traditionally, term generation has been approached from a purely
syntactic perspective. This has more recently been the area of work 
of some teams using Large Language Models. But generating terms
just from the syntax is very difficult to guide heuristically.
Our work trying to understand the framework proposed by Laurent Lafforgue
that connects Topos Theory \cite{lawvere1971} with artificial intelligence, lead to the following
redefinition of the problem: leveraging the categorical
semantics of type theories to reframe autonomous theorem proving
as a structural graph exploration problem. By translating the 
syntax of Generalized Algebraic Theories (GATs) and the Calculus
of Constructions (CoC) into categorical structures ---namely
Contextual Categories and Doctrines of Constructions--- we obtain
term models that inherently posses a rich, geometric graph structure.
The soundness of this translation is guaranteed by the term model being 
the initial object of the category of models of the underlying type theory.
These result is known as the \textit{initiality conjecture} from Voevodsky \cite{voevodsky2017}.
It was first proven by Streicher in \cite{streicher} for CoC, and more recently 
for a full-featured Martin-L\"of type theory in September 2020 by Brunerie and Lumsdaine
in \cite{bruneire2020}.

Withing this categorical framework, the problem of finding a proof
translates to finding a specific section (a morphism) withing a relativized
contextual category. 
With these idea in mind of using categories and neural networks, we formalized 
an expansion algorithm 
---a deterministic process that exhaustively builds the finite subcategories
representing the search space of terms--- specifically designed to work together
with Graph Neural Networks (GNNs), and try to overcome these intrinsic severe exponential growth
inherent to any algorithm trying to exhaustivity solve the problem.

Graph Neural Networks are used to navigate and prune the categorical search space. 
Considering our finite subcategories that form
the known search space as graphs, they are the perfect inputs 
for message-passing neural networks. By training a GNN to classify
nodes in the context graph (deciding witch paths to expand, keep
or delete), we can effectively guide the expansion algorithm,
bypassing the combinatorial explosion and paving the way for efficient,
autonomous theorem proving.

\section*{Objectives}

In this work we first aim to provide a proof of the initiality conjecture for GATs, together with an extension of Streicher's
proof for CoC together with a GAT signature.
Then, once established the mathematical significance of objects and morphisms in the term model,
we then present the expansion algorithm, an algorithm to construct any object or morphism in a Contextual 
Category or Contextual Doctrine, and that is designed to work together with a graph neural network to efficiently
solve the problem in hand.
Then we implement the expansion algorithm for simply typed combinatory logic, as a simple example to give intuition
about how the expansion process works. 
Finally we give a proposal of a GNN-based model thought to optimize term generation.

\section*{Thesis Outline}
This thesis is structured to guide the reader from the abstract categorical foundations
to the practical implementation of AI-guided theorem proving:
\begin{itemize}
    \item \textbf{Chapter 1: Syntax and Semantics of Generalized Algebraic Theories.}
        We introduce Cartmell's Generalized Algebraic Theories and define their syntax and semantics, as 
        they were originally presented in \cite{cartmell1978} \cite{cartmell1986} and incorporating some notions
        presented by Streicher in \cite{streicher}.
        Then we provide a concept of model of a GAT $U$ together with a proof of the initiality of the term model $C(U)$.
        Finally, define the problem of type inhabitation in categorical terms, and present the foundational Expansion Algorithm.
    \item \textbf{Chapter 2: Calculus of Constructions.} We step up the complexity to the Calculus of 
        Constructions (CoC), introducing product types and the universe of propositions. 
        We present the definition of CoC together with the construction 
        of Doctrines of Constructions as models for CoC, given by Streicher in \cite{streicher}.
        We extend the definition of CoC to include a GAT signature, define the new term model and prove its initiality.
        Finally, we extend our categorical problem definition and 
        Expansion Algorithm to handle this richer syntax.
    \item \textbf{Chapter 3: Implementation of the Expansion Algorithm.} We bring the theory into 
        practice. We demonstrate the raw Expansion Algorithm generating terms of Simply Typed Combinatory 
        Logic ($\textbf{CL}_\to$) using \textsc{implica}, a custom Rust-based graph engine. 
        Finally, we outline the architecture for integrating Graph Neural Networks
        using a novel unsupervised "Conjecturer-Prover" training paradigm to optimize the expansion process.
\end{itemize}